% !TEX program = xelatex
% AI-effects 프로젝트 --- 분모가 있는 상대 활용도 RU1·RU2(paper 25)와 분모가 없는 점유율 지표 A·B(paper 27)의 순위 비교.
% 2026-10-24 재설계: 활용도 1·2(대화 총량 필요)를 상대 활용도 RU1·RU2로 바꾸고, 지표를 AEI 업무용 대화 비중에서 직접 만듦.
% 표 본문: results/table/32_summary.tex, 32_g1_pos.tex, 32_g1_neg.tex, 32_g2_pos.tex, 32_g2_neg.tex,
%          32_top_use1.tex, 32_top_use2.tex, 32_ksco1.tex (code/04_15_analysis_kor_use_vs_share.do가 생성)
% 전체 자료: results/table/32_kor_use_vs_share.xlsx
% 관련: paper 21(연계표 배분), paper 25(RU1·RU2), paper 27(점유율 지표 A~D), paper 29(로그 점유율 회귀), paper 31(A와 B)
% 컴파일: xelatex (kotex 필요)
\documentclass[11pt]{article}

\usepackage{fontspec}
\usepackage{kotex}
\usepackage[a4paper,margin=2.0cm]{geometry}
\usepackage{longtable}
\usepackage{booktabs}
\usepackage{array}
\usepackage{amsmath}
\usepackage{xcolor}
\usepackage{hyperref}
\usepackage{enumitem}

\hypersetup{colorlinks=true,linkcolor=blue!50!black,urlcolor=blue!50!black}
\newcolumntype{L}[1]{>{\raggedright\arraybackslash}p{#1}}
\newcolumntype{R}[1]{>{\raggedleft\arraybackslash}p{#1}}
\setlength{\LTleft}{0pt}
\setlength{\LTright}{0pt}
\renewcommand{\arraystretch}{1.15}
\setlength{\tabcolsep}{3.5pt}

\title{분모가 있는 상대 활용도와 분모가 없는 점유율 지표의 순위 비교\\[2mm]
\large --- RU1 대 지표 A, RU2 대 지표 B ---}
\author{AI-effects 프로젝트}
\date{\today}

\begin{document}
\maketitle

\begin{abstract}
\noindent paper 25의 상대 활용도 RU1(대화 점유율 $\div$ 취업자 비중)과 RU2(절감 시간 점유율 $\div$ 근로시간 비중)는 분모에 직업 규모가 있다. paper 27의 지표 A(대화 점유율)와 B(절감 시간 점유율)는 분모가 없다. 정의상 $\ln \text{A} = \ln \mathrm{RU1} + \ln e$, $\ln \text{B} = \ln \mathrm{RU2} + \ln \eta$다($e$ = 취업자 비중, $\eta$ = 총 주업 근로시간 비중). 따라서 두 순위의 차이는 모두 직업 규모에서 나온다. KSCO 8차 소분류 129개에서 순위 격차 $G_1 = r(\mathrm{RU1}) - r(\text{A})$와 $G_2 = r(\mathrm{RU2}) - r(\text{B})$를 구했다(양수 = 점유율 순위가 더 높음). 2026년 10월 재설계로 활용도 1$\cdot$2(대화 총량이 필요한 수준값)를 RU1$\cdot$RU2로 바꿨다. RU는 예전 활용도를 전국값으로 나눈 것이라 순위가 거의 같고, 이 노트의 결론도 그대로다.
\begin{itemize}[leftmargin=*]
  \item \textbf{순위는 꽤 다르다.} Spearman 순위상관은 RU1--A 0.81, RU2--B 0.86이다. 상위 10 중 겹치는 직업은 6개, 5개, 상위 30 중에서는 각각 19개다. 격차의 절댓값은 평균 17.6계단, 14.8계단이고, 20계단 이상인 소분류가 44개, 39개다. 같은 계열끼리의 순위상관(RU1--RU2 0.98, A--B 0.96)보다 훨씬 낮다.
  \item \textbf{대화 점유율은 직업 규모와 무관하다.} $\ln \text{A}$를 $\ln E$에 회귀한 기울기는 $-0.01$($R^2$ = 0.00), $\ln \text{B}$를 $\ln H$에 회귀한 기울기는 $-0.16$($R^2$ = 0.01)다. 취업자가 많은 직업이라고 AI 사용 몫이 더 크지 않다. 그래서 규모로 나눈 RU는 규모와 음의 상관을 갖는다(RU1--취업자 순위상관 $-0.55$, RU2 $-0.54$). 반면 A, B와 취업자의 상관은 $-0.01$, $-0.10$으로 거의 0이다.
  \item \textbf{격차는 사실상 직업 규모의 순위다.} $G_1$과 취업자 수의 순위상관은 0.89, $G_2$와 총 근로시간은 0.87이다. RU에서는 낮지만 점유율에서 높은 직업은 큰 사무$\cdot$판매 직업이다. 521 매장 판매($G_1$ = +65), 331 회계 및 경리(+49), 321 자재 및 생산관리(+43), 311 기획 및 마케팅(+41) 사무원이 그렇다. 반대로 RU에서 높지만 점유율에서 낮은 직업은 작은 관리자$\cdot$장치 직업이다. 159 기타 판매 및 고객 서비스 관리자($-74$), 134 문화$\cdot$예술 관련 관리자($-64$), 226 통신 및 방송 송출 장비 기사($-50$), 882 재활용 처리 및 소각로 조작원($-44$), 149 기타 건설$\cdot$전기 및 생산 관련 관리자($-43$)가 그렇다.
  \item \textbf{RU 상위에는 표본이 작은 직업이 섞인다.} 지역별고용조사 표본 30 미만인 소분류 6개 중 2개가 RU1 상위 10에, 3개가 RU2 상위 10에 든다. A, B 상위 10에는 하나도 없다. 149(RU1 3위, 표본 9)는 A로는 46위다.
  \item \textbf{대분류}: 관리자는 RU1 상위 30에 10개가 들지만 A 상위 30에는 3개다(평균 $G_1$ $-26.8$). 사무 종사자는 1개에서 4개로 는다(평균 $G_1$ +24.7).
\end{itemize}
\end{abstract}

\section{두 지표의 관계}

소분류 $j$의 지표 A$\cdot$B(\%)와 2022년 상반기 취업자 수 $E_j$, 월 총 주업 근로시간 $H_j$로
\[
\mathrm{RU1}_j = \frac{\text{A}_j}{e_j}, \quad \mathrm{RU2}_j = \frac{\text{B}_j}{\eta_j}, \quad
e_j = 100\,\frac{E_j}{\sum_k E_k}, \quad \eta_j = 100\,\frac{H_j}{\sum_k H_k}
\]
이다(paper 25, 27과 같음; 분모 합은 소분류 163개 전체). 그러므로
\[
\ln \text{A}_j = \ln \mathrm{RU1}_j + \ln e_j, \qquad
\ln \text{B}_j = \ln \mathrm{RU2}_j + \ln \eta_j
\]
가 항등식으로 성립한다(do 파일에서 assert). $\ln e_j = \ln E_j + \text{상수}$이므로 아래 분해와 회귀에서는 $\ln E$, $\ln H$를 그대로 쓴다. RU는 ``규모 대비 얼마나 쓰는가''(강도), 점유율은 ``전체 사용 중 얼마가 이 직업에 있는가''(규모를 포함한 양)이다. 두 순위가 다른 이유는 규모 $E_j$, $H_j$ 하나뿐이다. $H_j = E_j \times$ 1인당 근로시간이므로 RU2--B 쌍에는 1인당 근로시간의 차이도 들어 있다(소분류 주당 22$\sim$47시간).

순위 격차는
\[
G_{1j} = r(\mathrm{RU1}_j) - r(\text{A}_j), \qquad G_{2j} = r(\mathrm{RU2}_j) - r(\text{B}_j)
\]
로 쟀다. $r(\cdot)$은 129개 안 순위(1 = 최상위, 같은 값은 소분류 코드 순)다. 양수이면 점유율 순위가 RU 순위보다 높은 직업, 즉 규모가 큰 직업이다.

규모가 얼마나 순위를 바꾸는지는 분산 분해와 규모 탄력성으로 본다.
\[
\mathrm{var}(\ln \text{A}) = \mathrm{var}(\ln \mathrm{RU1}) + \mathrm{var}(\ln E) + 2\,\mathrm{cov}(\ln \mathrm{RU1}, \ln E)
\]
이고, $\ln \text{A}$를 $\ln E$에 회귀한 기울기는 $1 + \mathrm{cov}(\ln \mathrm{RU1}, \ln E)/\mathrm{var}(\ln E)$다. 이 기울기(규모 탄력성)가 1이면 사용이 취업자에 비례해 RU가 규모와 무관하고, 0이면 사용이 규모와 무관해 RU가 $1/E$처럼 움직인다.

\section{결과}

\subsection{요약}

\begin{longtable}{L{12.5cm}R{2.5cm}}
\caption{상대 활용도와 점유율 지표의 순위 관계, 분산 분해, 규모 탄력성(소분류 129개)}\label{tab:summary}\\
\toprule
항목 & 값 \\
\midrule
\endfirsthead
\toprule
항목 & 값 \\
\midrule
\endhead
% 04_analysis_kor_use_vs_share.do가 생성. 직접 고치지 말 것.
활용도 1과 지표 A의 Spearman 순위상관 & 0.80 \\
활용도 2와 지표 B의 Spearman 순위상관 & 0.86 \\
(참고) 활용도 1과 활용도 2 & 0.97 \\
(참고) 지표 A와 지표 B & 0.96 \\
상위 10 겹침: 활용도 1 -- A & 6 \\
상위 30 겹침: 활용도 1 -- A & 19 \\
상위 10 겹침: 활용도 2 -- B & 5 \\
상위 30 겹침: 활용도 2 -- B & 19 \\
\(|G_1|\) 평균 & 18.0 \\
\(|G_1|\) 최댓값 & 74 \\
\(|G_1| \ge 20\)인 소분류 수 & 48 \\
\(|G_2|\) 평균 & 15.0 \\
\(|G_2|\) 최댓값 & 60 \\
\(|G_2| \ge 20\)인 소분류 수 & 39 \\
\(G_1\)과 취업자 수의 Spearman 순위상관 & 0.90 \\
\(G_2\)와 총 주업 근로시간의 Spearman 순위상관 & 0.87 \\
\(G_1\)과 \(G_2\)의 Spearman 순위상관 & 0.96 \\
\(\mathrm{var}(\ln \mathrm{A})\) & 4.05 \\
\quad 활용도 1 분산의 몫 & 1.52 \\
\quad 취업자 수 분산의 몫 & 0.52 \\
\quad 공분산(\(2\,\mathrm{cov}\))의 몫 & -1.04 \\
\(\ln\) 활용도 1과 \(\ln E\)의 Pearson 상관 & -0.59 \\
\(\ln \mathrm{A}\)를 \(\ln E\)에 회귀한 기울기(규모 탄력성) & 0.00 \\
\quad 같은 회귀의 \(R^2\) & 0.00 \\
\(\mathrm{var}(\ln \mathrm{B})\) & 5.97 \\
\quad 활용도 2 분산의 몫 & 1.45 \\
\quad 총 주업 근로시간 분산의 몫 & 0.35 \\
\quad 공분산(\(2\,\mathrm{cov}\))의 몫 & -0.80 \\
\(\ln\) 활용도 2와 \(\ln H\)의 Pearson 상관 & -0.56 \\
\(\ln \mathrm{B}\)를 \(\ln H\)에 회귀한 기울기(규모 탄력성) & -0.14 \\
\quad 같은 회귀의 \(R^2\) & 0.01 \\
표본 30 미만 소분류 수: 활용도 1 상위 10 중 & 2 \\
표본 30 미만 소분류 수: 지표 A 상위 10 중 & 0 \\
표본 30 미만 소분류 수: 활용도 2 상위 10 중 & 3 \\
표본 30 미만 소분류 수: 지표 B 상위 10 중 & 0 \\
표본 30 미만 소분류 수(129개 중) & 6 \\
취업자 수 최댓값 / 최솟값(배) & 1,282 \\
활용도 1과 취업자 수의 Spearman 순위상관 & -0.56 \\
지표 A와 취업자 수의 Spearman 순위상관 & -0.01 \\
활용도 2와 취업자 수의 Spearman 순위상관 & -0.54 \\
지표 B와 취업자 수의 Spearman 순위상관 & -0.10 \\

\bottomrule
\end{longtable}

\paragraph{규모 탄력성이 0이다.} $\ln \text{A}$를 $\ln E$에 회귀한 기울기는 $-0.01$이고 $R^2$는 0.00이다. 취업자 수가 1,282배 차이 나는 129개 소분류에서 업무용 대화 점유율은 취업자 수와 아무 관계가 없다. 그래서 $\ln \mathrm{RU1} = \ln \text{A} - \ln e$의 분산은 $\ln \text{A}$의 분산에 $\ln E$의 분산을 더한 것과 거의 같다(1.50 $\approx$ 1 + 0.49). RU1과 취업자 수의 상관($-0.58$, 로그 Pearson)은 순전히 분모에서 생긴다. RU2와 B도 같다(기울기 $-0.16$, $R^2$ = 0.01). 이 결과는 paper 29에서 로그 점유율을 노출도와 로그 취업자 수에 회귀했을 때 로그 취업자 수의 계수가 0과 구별되지 않았던 것과 같은 사실이다. 재설계 뒤 $\mathrm{var}(\ln \text{A})$가 4.05에서 4.34로 커진 것은 한 달만 공개된 하위권 소분류 7개의 점유율이 절반이 되어 아래 꼬리가 길어졌기 때문이다(paper 25).

이 ``규모 탄력성 0''에는 두 가지가 섞여 있다.
\begin{enumerate}[leftmargin=*]
  \item \textbf{구축 방식.} 소분류 점유율은 AEI의 SOC별 한국 업무용 대화 비중을 SOC--KSCO 연계표로 나눠 만든다. 배분 가중치 \texttt{w\_alloc}은 단계별 \emph{균등} 배분이다(paper 21). 한 SOC가 여러 소분류에 이어지면 취업자가 많은 소분류도 적은 소분류도 같은 몫을 받는다. 따라서 사용이 한국 취업자 규모를 따라가게 만드는 장치가 자료에 없다.
  \item \textbf{AEI의 직업 정의.} AEI의 SOC는 사용자의 직업이 아니라 대화 과업이 속한 직업이다. 매장 판매원이 많아도 ``Cashiers'' 과업의 대화가 많아지는 것은 아니다.
\end{enumerate}
둘 중 무엇이 더 큰지는 이 자료로 가를 수 없다. 다만 어느 쪽이든 결과는 같다. \textbf{RU1$\cdot$RU2의 직업 간 차이 중 상당 부분은 분모(직업 규모)의 차이}이고, 취업자가 적은 직업일수록 RU가 기계적으로 커진다(paper 27이 점유율 지표를 만든 이유).

\paragraph{순위의 일치.} 순위상관은 RU1--A 0.81, RU2--B 0.86이다. RU2--B가 조금 높은 이유는 $\ln H$의 분산(var($\ln$ B)의 0.32)이 $\ln E$의 분산(var($\ln$ A)의 0.49)보다 상대적으로 작기 때문이다. 상위 10 겹침은 6개, 5개, 상위 30 겹침은 각각 19개다. 격차의 절댓값이 20계단 이상인 소분류가 44개, 39개로, paper 31의 A--B 비교(9개)보다 훨씬 많다. 즉 ``대화냐 절감 시간이냐''보다 ``분모가 있느냐 없느냐''가 순위를 훨씬 크게 바꾼다. $G_1$과 $G_2$의 순위상관은 0.96이라, 두 쌍에서 순위가 바뀌는 직업은 거의 같다.

\subsection{격차가 큰 소분류}

표의 $\dagger$는 지역별고용조사 2022년 상반기 표본이 30 미만인 소분류다(분모가 불안정).

\begin{small}
\begin{longtable}{R{0.5cm}L{4.4cm}R{1.3cm}R{1.2cm}R{1.0cm}R{0.9cm}R{0.9cm}R{1.4cm}R{1.2cm}}
\caption{RU1 대 지표 A: 점유율 순위가 더 높은 소분류($G_1$ 상위 15)}\label{tab:g1pos}\\
\toprule
 & 소분류 & RU1 & $r$(RU1) & A(\%) & $r$(A) & $G_1$ & 취업자(천 명) & $r$(취업자) \\
\midrule
\endfirsthead
% 04_analysis_kor_use_vs_share.do가 생성. 직접 고치지 말 것.
1 & 521 매장 판매 종사자 & 1.71 & 95 & 0.83 & 30 & +65 & 1,683 & 1 \\
2 & 331 회계 및 경리 사무원 & 1.85 & 94 & 0.44 & 45 & +49 & 826 & 5 \\
3 & 321 자재 및 생산관리 사무원 & 1.33 & 98 & 0.29 & 54 & +44 & 750 & 7 \\
4 & 311 기획 및 마케팅 사무원 & 1.01 & 106 & 0.21 & 63 & +43 & 707 & 8 \\
5 & 352 정부 행정 사무 종사자 & 4.16 & 77 & 0.63 & 36 & +41 & 523 & 12 \\
6 & 371 총무 사무원 및 대학 행정조교 & 7.80 & 59 & 1.33 & 18 & +41 & 591 & 10 \\
7 & 510 영업 종사자 & 5.23 & 69 & 0.92 & 28 & +41 & 609 & 9 \\
8 & 264 문리‧기술 및 예능 강사 & 10.97 & 54 & 1.85 & 14 & +40 & 585 & 11 \\
9 & 295 디자이너 & 15.47 & 51 & 1.19 & 20 & +31 & 266 & 28 \\
10 & 340 금융 사무 종사자 & 31.50 & 41 & 2.61 & 10 & +31 & 287 & 25 \\
11 & 372 비서 및 사무 지원 종사자 & 46.22 & 33 & 4.69 & 4 & +29 & 351 & 21 \\
12 & 284 감정‧기술 영업 및 중개 관련 종사자 & 9.66 & 55 & 1.01 & 27 & +28 & 361 & 20 \\
13 & 262 학교 교사 & 1.18 & 101 & 0.14 & 74 & +27 & 422 & 16 \\
14 & 249 기타 보건 전문가 및 관련 종사자 & 3.68 & 79 & 0.30 & 53 & +26 & 280 & 26 \\
15 & 910 건설 및 광업 단순 종사자 & 0.58 & 115 & 0.07 & 91 & +24 & 433 & 15 \\

\bottomrule
\end{longtable}

\begin{longtable}{R{0.5cm}L{4.4cm}R{1.3cm}R{1.2cm}R{1.0cm}R{0.9cm}R{0.9cm}R{1.4cm}R{1.2cm}}
\caption{RU1 대 지표 A: RU 순위가 더 높은 소분류($G_1$ 하위 15)}\label{tab:g1neg}\\
\toprule
 & 소분류 & RU1 & $r$(RU1) & A(\%) & $r$(A) & $G_1$ & 취업자(천 명) & $r$(취업자) \\
\midrule
\endfirsthead
% 04_15_analysis_kor_use_vs_share.do가 생성. 직접 고치지 말 것.
1 & 159 기타 판매 및 고객 서비스 관리자$^\dagger$ & 18.21 & 12 & 0.09 & 86 & -74 & 1 & 129 \\
2 & 134 문화‧예술 관련 관리자$^\dagger$ & 6.81 & 25 & 0.08 & 89 & -64 & 3 & 127 \\
3 & 226 통신 및 방송 송출 장비 기사 & 3.61 & 34 & 0.10 & 84 & -50 & 8 & 122 \\
4 & 882 재활용 처리 및 소각로 조작원 & 0.60 & 64 & 0.02 & 108 & -44 & 11 & 118 \\
5 & 149 기타 건설‧전기 및 생산 관련 관리자$^\dagger$ & 46.03 & 3 & 0.42 & 46 & -43 & 3 & 128 \\
6 & 153 환경‧청소 및 경비 관련 관리자$^\dagger$ & 17.27 & 13 & 0.27 & 55 & -42 & 4 & 125 \\
7 & 441 운송 서비스 종사자 & 0.61 & 63 & 0.04 & 101 & -38 & 17 & 109 \\
8 & 881 상하수도 처리 장치 조작원 & 0.42 & 70 & 0.02 & 107 & -37 & 16 & 110 \\
9 & 112 기업 대표 및 기업 고위 임원 & 3.16 & 37 & 0.15 & 72 & -35 & 13 & 113 \\
10 & 791 공예 및 귀금속 세공원 & 0.68 & 57 & 0.07 & 92 & -35 & 29 & 99 \\
11 & 111 의회 의원‧고위 공무원 및 공공단체 임원 & 6.87 & 24 & 0.25 & 58 & -34 & 10 & 119 \\
12 & 233 금속‧재료공학 기술자 및 시험원 & 4.99 & 28 & 0.22 & 62 & -34 & 12 & 115 \\
13 & 432 혼례 및 장례 종사자 & 2.10 & 44 & 0.13 & 77 & -33 & 17 & 108 \\
14 & 152 고객 서비스 관리자 & 3.80 & 32 & 0.20 & 64 & -32 & 15 & 112 \\
15 & 252 시민사회 활동가$^\dagger$ & 0.16 & 93 & 0.00 & 125 & -32 & 4 & 126 \\

\bottomrule
\end{longtable}
\end{small}

\paragraph{점유율 순위가 높은 직업($G_1 > 0$).} 상위 15개 모두 취업자 순위 30위 안이다. 521 매장 판매 종사자는 취업자 1위(168만 명)다. RU1은 0.14(규모의 7분의 1)로 95위지만, 사용 몫 자체는 커서 A로는 30위다. 331 회계 및 경리(취업자 5위), 321 자재 및 생산관리(7위), 311 기획 및 마케팅(8위), 352 정부 행정 사무(12위), 371 총무 사무원(10위), 510 영업(9위), 264 문리$\cdot$기술 및 예능 강사(11위) 등 큰 사무$\cdot$영업$\cdot$교육 직업이 이어진다. 이 직업들은 규모 대비 사용은 많지 않다. 하지만 사람이 많아 AI 사용 총량에서는 무시할 수 없는 몫을 차지한다. 일자리 영향 시나리오(paper 27)에서 RU로 몫을 나누면 이 직업들의 몫은 작게 잡힌다.

\paragraph{RU 순위가 높은 직업($G_1 < 0$).} 상위 15개 모두 취업자 순위 99위 이하다. 159 기타 판매 및 고객 서비스 관리자는 취업자가 약 1천 명(표본 6)이라 RU1이 18.2로 12위지만, A로는 86위다. 134 문화$\cdot$예술 관련 관리자, 149 기타 건설$\cdot$전기 및 생산 관련 관리자, 153 환경$\cdot$청소 및 경비 관련 관리자, 252 시민사회 활동가처럼 표본 30 미만($\dagger$)인 소분류가 5개 들어 있다. RU가 높은 것이 실제 사용 강도인지, 작은 분모에서 생긴 값인지 가르기 어려운 직업들이다.

\begin{small}
\begin{longtable}{R{0.5cm}L{4.2cm}R{1.3cm}R{1.2cm}R{1.0cm}R{0.9cm}R{0.9cm}R{1.3cm}R{1.0cm}R{1.0cm}}
\caption{RU2 대 지표 B: 점유율 순위가 더 높은 소분류($G_2$ 상위 15)}\label{tab:g2pos}\\
\toprule
 & 소분류 & RU2 & $r$(RU2) & B(\%) & $r$(B) & $G_2$ & 취업자(천 명) & 주당 시간 & $r(H)$ \\
\midrule
\endfirsthead
% 04_15_analysis_kor_use_vs_share.do가 생성. 직접 고치지 말 것.
1 & 521 매장 판매 종사자 & 0.04 & 105 & 0.26 & 55 & +50 & 1,683 & 41.4 & 1 \\
2 & 264 문리‧기술 및 예능 강사 & 1.32 & 50 & 2.07 & 13 & +37 & 585 & 29.2 & 14 \\
3 & 371 총무 사무원 및 대학 행정조교 & 0.47 & 61 & 1.01 & 26 & +35 & 591 & 39.2 & 10 \\
4 & 321 자재 및 생산관리 사무원 & 0.07 & 94 & 0.21 & 60 & +34 & 750 & 41.7 & 6 \\
5 & 510 영업 종사자 & 0.23 & 72 & 0.50 & 39 & +33 & 609 & 39.7 & 9 \\
6 & 340 금융 사무 종사자 & 1.70 & 48 & 1.78 & 16 & +32 & 287 & 39.8 & 24 \\
7 & 331 회계 및 경리 사무원 & 0.08 & 90 & 0.23 & 59 & +31 & 826 & 38.9 & 5 \\
8 & 352 정부 행정 사무 종사자 & 0.24 & 71 & 0.45 & 40 & +31 & 523 & 39.5 & 11 \\
9 & 311 기획 및 마케팅 사무원 & 0.02 & 111 & 0.06 & 81 & +30 & 707 & 40.1 & 8 \\
10 & 941 청소 관련 종사자 & 0.02 & 114 & 0.05 & 84 & +30 & 1,141 & 26.0 & 7 \\
11 & 284 감정‧기술 영업 및 중개 관련 종사자 & 0.54 & 58 & 0.75 & 31 & +27 & 361 & 41.4 & 18 \\
12 & 295 디자이너 & 1.05 & 52 & 1.01 & 25 & +27 & 266 & 39.4 & 28 \\
13 & 372 비서 및 사무 지원 종사자 & 4.65 & 31 & 5.24 & 4 & +27 & 351 & 34.9 & 21 \\
14 & 234 전기‧전자공학 기술자 및 시험원 & 0.95 & 54 & 0.83 & 29 & +25 & 228 & 41.7 & 31 \\
15 & 222 컴퓨터 시스템 및 소프트웨어 전문가 & 8.15 & 22 & 12.21 & 1 & +21 & 399 & 40.8 & 15 \\

\bottomrule
\end{longtable}

\begin{longtable}{R{0.5cm}L{4.2cm}R{1.3cm}R{1.2cm}R{1.0cm}R{0.9cm}R{0.9cm}R{1.3cm}R{1.0cm}R{1.0cm}}
\caption{RU2 대 지표 B: RU 순위가 더 높은 소분류($G_2$ 하위 15)}\label{tab:g2neg}\\
\toprule
 & 소분류 & RU2 & $r$(RU2) & B(\%) & $r$(B) & $G_2$ & 취업자(천 명) & 주당 시간 & $r(H)$ \\
\midrule
\endfirsthead
% 04_15_analysis_kor_use_vs_share.do가 생성. 직접 고치지 말 것.
1 & 159 기타 판매 및 고객 서비스 관리자$^\dagger$ & 24.55 & 8 & 0.13 & 68 & -60 & 1 & 42.4 & 129 \\
2 & 134 문화‧예술 관련 관리자$^\dagger$ & 7.35 & 23 & 0.08 & 78 & -55 & 3 & 37.9 & 126 \\
3 & 882 재활용 처리 및 소각로 조작원 & 0.57 & 57 & 0.02 & 101 & -44 & 11 & 41.0 & 118 \\
4 & 252 시민사회 활동가$^\dagger$ & 0.14 & 79 & 0.00 & 122 & -43 & 4 & 30.1 & 127 \\
5 & 153 환경‧청소 및 경비 관련 관리자$^\dagger$ & 19.26 & 12 & 0.33 & 50 & -38 & 4 & 42.4 & 125 \\
6 & 135 정보 통신 관련 관리자 & 10.09 & 20 & 0.24 & 57 & -37 & 6 & 41.2 & 123 \\
7 & 226 통신 및 방송 송출 장비 기사 & 5.56 & 29 & 0.17 & 63 & -34 & 8 & 42.2 & 122 \\
8 & 881 상하수도 처리 장치 조작원 & 0.39 & 66 & 0.02 & 100 & -34 & 16 & 42.1 & 108 \\
9 & 149 기타 건설‧전기 및 생산 관련 관리자$^\dagger$ & 68.95 & 3 & 0.65 & 34 & -31 & 3 & 39.9 & 128 \\
10 & 432 혼례 및 장례 종사자 & 1.38 & 49 & 0.08 & 77 & -28 & 17 & 38.3 & 109 \\
11 & 133 보건 및 사회복지 관련 관리자 & 7.04 & 25 & 0.31 & 52 & -27 & 11 & 42.3 & 116 \\
12 & 111 의회 의원‧고위 공무원 및 공공단체 임원 & 9.72 & 21 & 0.37 & 47 & -26 & 10 & 41.0 & 119 \\
13 & 112 기업 대표 및 기업 고위 임원 & 4.46 & 32 & 0.23 & 58 & -26 & 13 & 41.9 & 112 \\
14 & 441 운송 서비스 종사자 & 0.09 & 86 & 0.00 & 112 & -26 & 17 & 32.5 & 114 \\
15 & 233 금속‧재료공학 기술자 및 시험원 & 7.03 & 26 & 0.32 & 51 & -25 & 12 & 41.1 & 115 \\

\bottomrule
\end{longtable}
\end{small}

\paragraph{RU2 대 B.} 목록은 RU1 대 A와 거의 같다. 521 매장 판매($G_2$ = +50), 사무 직업들, 그리고 159$\cdot$134$\cdot$252$\cdot$153$\cdot$149 같은 작은 관리자$\cdot$단체 직업이 양 끝에 온다. 다른 점은 두 가지다. 첫째, 규모가 취업자 수가 아니라 총 근로시간이라 1인당 근로시간이 순위에 들어간다. 941 청소 관련 종사자는 취업자로는 3위지만 주당 26시간이라 총 근로시간으로는 7위다. 둘째, 222 컴퓨터 시스템 및 소프트웨어 전문가는 RU2로는 22위지만 B로는 1위다(+21). 취업자가 40만 명이고 대화당 절감 시간도 길기 때문이다(paper 31).

\subsection{RU 상위 20의 점유율 순위}

\begin{small}
\begin{longtable}{R{0.6cm}L{4.6cm}R{1.3cm}R{1.0cm}R{0.9cm}R{0.9cm}R{1.4cm}R{1.4cm}R{1.2cm}}
\caption{RU1 상위 20 소분류의 지표 A 순위(표본 = 지역별고용조사 2022 상반기 표본 수)}\label{tab:topuse1}\\
\toprule
$r$(RU1) & 소분류 & RU1 & A(\%) & $r$(A) & $G_1$ & 취업자(천 명) & 취업자 비중(\%) & 표본 \\
\midrule
\endfirsthead
% 04_analysis_kor_use_vs_share.do가 생성. 직접 고치지 말 것.
1 & 213 생명 및 자연과학 관련 시험원 & 839.83 & 2.79 & 8 & -7 & 11 & 0.05 & 84 \\
2 & 332 통계‧데이터 관련 사무원 & 569.73 & 2.61 & 9 & -7 & 16 & 0.06 & 127 \\
3 & 149 기타 건설‧전기 및 생산 관련 관리자$^\dagger$ & 567.32 & 0.42 & 46 & -43 & 3 & 0.01 & 9 \\
4 & 291 작가 및 언론 관련 전문가 & 430.04 & 11.92 & 1 & +3 & 96 & 0.39 & 564 \\
5 & 292 학예사‧사서 및 기록물 관리사 & 422.60 & 3.81 & 7 & -2 & 31 & 0.13 & 228 \\
6 & 122 마케팅 및 광고‧홍보 관리자$^\dagger$ & 399.47 & 1.01 & 25 & -19 & 9 & 0.04 & 28 \\
7 & 223 네트워크 및 정보 보안 전문가 & 397.80 & 3.84 & 6 & +1 & 33 & 0.14 & 152 \\
8 & 225 정보 시스템 및 웹 운영자 & 365.21 & 8.31 & 3 & +5 & 79 & 0.32 & 421 \\
9 & 135 정보 통신 관련 관리자 & 296.97 & 0.54 & 39 & -30 & 6 & 0.03 & 35 \\
10 & 298 문화‧예술 관련 기획자 & 234.81 & 0.66 & 35 & -25 & 10 & 0.04 & 64 \\
11 & 283 상품 기획‧홍보 및 조사 전문가 & 233.70 & 4.20 & 5 & +6 & 62 & 0.25 & 345 \\
12 & 159 기타 판매 및 고객 서비스 관리자$^\dagger$ & 224.40 & 0.09 & 86 & -74 & 1 & 0.01 & 6 \\
13 & 153 환경‧청소 및 경비 관련 관리자$^\dagger$ & 212.90 & 0.27 & 55 & -42 & 4 & 0.02 & 20 \\
14 & 212 인문 및 사회과학 전문가 & 200.77 & 1.26 & 19 & -5 & 22 & 0.09 & 136 \\
15 & 224 데이터 전문가 & 152.15 & 1.10 & 21 & -6 & 25 & 0.10 & 113 \\
16 & 282 금융 및 보험 전문가 & 145.59 & 1.98 & 13 & +3 & 47 & 0.19 & 289 \\
17 & 522 상품 대여 종사자 & 122.18 & 1.07 & 22 & -5 & 30 & 0.12 & 220 \\
18 & 852 냉난방 관련 설비 조작원 & 116.84 & 0.59 & 37 & -19 & 17 & 0.07 & 135 \\
19 & 271 법률 전문가 & 109.45 & 1.04 & 24 & -5 & 33 & 0.13 & 212 \\
20 & 269 기타 교육 전문가 & 106.87 & 1.41 & 16 & +4 & 46 & 0.19 & 277 \\

\bottomrule
\end{longtable}

\begin{longtable}{R{0.6cm}L{4.6cm}R{1.3cm}R{1.0cm}R{0.9cm}R{0.9cm}R{1.4cm}R{1.4cm}R{1.2cm}}
\caption{RU2 상위 20 소분류의 지표 B 순위}\label{tab:topuse2}\\
\toprule
$r$(RU2) & 소분류 & RU2 & B(\%) & $r$(B) & $G_2$ & 취업자(천 명) & 취업자 비중(\%) & 표본 \\
\midrule
\endfirsthead
% 04_15_analysis_kor_use_vs_share.do가 생성. 직접 고치지 말 것.
1 & 213 생명 및 자연과학 관련 시험원 & 85.97 & 3.62 & 8 & -7 & 11 & 0.05 & 84 \\
2 & 332 통계‧데이터 관련 사무원 & 74.38 & 4.03 & 6 & -4 & 16 & 0.06 & 127 \\
3 & 149 기타 건설‧전기 및 생산 관련 관리자$^\dagger$ & 68.95 & 0.65 & 34 & -31 & 3 & 0.01 & 9 \\
4 & 122 마케팅 및 광고‧홍보 관리자$^\dagger$ & 46.32 & 1.50 & 18 & -14 & 9 & 0.04 & 28 \\
5 & 291 작가 및 언론 관련 전문가 & 32.51 & 10.47 & 2 & +3 & 96 & 0.39 & 564 \\
6 & 223 네트워크 및 정보 보안 전문가 & 29.55 & 3.74 & 7 & -1 & 33 & 0.14 & 152 \\
7 & 298 문화‧예술 관련 기획자 & 27.29 & 0.98 & 27 & -20 & 10 & 0.04 & 64 \\
8 & 159 기타 판매 및 고객 서비스 관리자$^\dagger$ & 24.55 & 0.13 & 68 & -60 & 1 & 0.01 & 6 \\
9 & 283 상품 기획‧홍보 및 조사 전문가 & 22.89 & 5.15 & 5 & +4 & 62 & 0.25 & 345 \\
10 & 292 학예사‧사서 및 기록물 관리사 & 21.37 & 2.29 & 12 & -2 & 31 & 0.13 & 228 \\
11 & 224 데이터 전문가 & 19.34 & 1.82 & 15 & -4 & 25 & 0.10 & 113 \\
12 & 153 환경‧청소 및 경비 관련 관리자$^\dagger$ & 19.26 & 0.33 & 50 & -38 & 4 & 0.02 & 20 \\
13 & 225 정보 시스템 및 웹 운영자 & 18.33 & 5.34 & 3 & +10 & 79 & 0.32 & 421 \\
14 & 212 인문 및 사회과학 전문가 & 14.86 & 1.12 & 23 & -9 & 22 & 0.09 & 136 \\
15 & 269 기타 교육 전문가 & 14.33 & 2.02 & 14 & +1 & 46 & 0.19 & 277 \\
16 & 282 금융 및 보험 전문가 & 13.74 & 2.48 & 11 & +5 & 47 & 0.19 & 289 \\
17 & 211 생명 및 자연과학 관련 전문가 & 10.30 & 2.89 & 9 & +8 & 76 & 0.31 & 442 \\
18 & 271 법률 전문가 & 10.21 & 1.25 & 22 & -4 & 33 & 0.13 & 212 \\
19 & 852 냉난방 관련 설비 조작원 & 10.15 & 0.70 & 33 & -14 & 17 & 0.07 & 135 \\
20 & 135 정보 통신 관련 관리자 & 10.09 & 0.24 & 57 & -37 & 6 & 0.03 & 35 \\

\bottomrule
\end{longtable}
\end{small}

RU1 상위 10 중 A 상위 10에도 드는 직업은 213 생명 및 자연과학 관련 시험원(A 8위), 332 통계$\cdot$데이터 관련 사무원(9위), 291 작가 및 언론(1위), 292 학예사$\cdot$사서 및 기록물 관리사(7위), 223 네트워크 및 정보 보안(6위), 225 정보 시스템 및 웹 운영자(3위)다. 이들은 규모 대비 사용도 많고 총량도 커서 어느 지표로 보든 AI 사용의 중심이다. 빠지는 4개는 149(A 46위$\dagger$), 122 마케팅 및 광고$\cdot$홍보 관리자(25위$\dagger$), 135 정보 통신 관련 관리자(39위), 298 문화$\cdot$예술 관련 기획자(35위)로, 모두 취업자 비중 0.04\% 이하인 작은 직업이다. RU2 상위 10에서도 같은 구조가 반복되며, 159(B 68위$\dagger$)가 8위로 들어온다.

\subsection{대분류}

\begin{longtable}{L{4.2cm}R{0.9cm}R{1.3cm}R{1.2cm}R{1.0cm}R{1.2cm}R{1.0cm}R{1.0cm}R{1.0cm}R{1.2cm}R{1.2cm}}
\caption{대분류별 상위 30 소분류 수, 점유율 합, 평균 순위 격차(격차는 소분류 단순 평균)}\label{tab:ksco1}\\
\toprule
대분류 & 소분류 & 취업자(\%) & RU1 상위 30 & A 상위 30 & RU2 상위 30 & B 상위 30 & A(\%) & B(\%) & 평균 $G_1$ & 평균 $G_2$ \\
\midrule
\endfirsthead
% 04_15_analysis_kor_use_vs_share.do가 생성. 직접 고치지 말 것.
1 관리자 & 16 & 1.7 & 10 & 3 & 11 & 4 & 6.9 & 8.9 & -26.8 & -19.7 \\
2 전문가 및 관련 종사자 & 45 & 23.6 & 17 & 20 & 17 & 22 & 70.5 & 72.2 & +2.8 & +2.8 \\
3 사무 종사자 & 13 & 19.8 & 1 & 4 & 1 & 4 & 14.6 & 14.4 & +24.7 & +19.8 \\
4 서비스 종사자 & 13 & 11.4 & 0 & 0 & 0 & 0 & 1.1 & 0.5 & -2.2 & -4.3 \\
5 판매 종사자 & 5 & 10.8 & 1 & 3 & 0 & 0 & 3.2 & 1.3 & +22.6 & +15.8 \\
6 농림∙어업 숙련 종사자 & 4 & 5.9 & 0 & 0 & 0 & 0 & 0.0 & 0.0 & -2.8 & -1.0 \\
7 기능원 및 관련 기능 종사자 & 14 & 7.0 & 0 & 0 & 0 & 0 & 1.9 & 1.4 & -2.7 & -2.1 \\
8 장치∙기계 조작 및 조립 종사자 & 10 & 6.6 & 1 & 0 & 1 & 0 & 0.9 & 1.0 & -12.7 & -11.1 \\
9 단순노무 종사자 & 9 & 13.1 & 0 & 0 & 0 & 0 & 0.9 & 0.4 & +8.1 & +5.8 \\

\bottomrule
\end{longtable}

관리자는 취업자의 1.7\%지만 RU1 상위 30에 10개, RU2 상위 30에 11개가 든다. 점유율로 보면 A 상위 30에 3개, B 상위 30에 4개다(평균 $G_1$ $-26.8$, $G_2$ $-19.7$). RU로 보면 ``관리자가 AI를 가장 많이 쓰는 직업군''으로 보이지만, 그 상당 부분은 작은 분모 때문이다. 사무 종사자(취업자 19.8\%)는 반대로 RU1 상위 30에 1개뿐이지만 A 상위 30에는 4개가 든다(평균 $G_1$ +24.7). 판매(+22.6)와 단순노무(+8.1)도 점유율 쪽 순위가 높다. 전문가 대분류는 어느 지표로 보든 상위 30의 절반 이상(17$\sim$22개)을 차지하고 평균 격차도 작다(+2.8, +2.8).

\section{해석과 유의사항}

\begin{enumerate}[leftmargin=*]
  \item \textbf{두 지표는 다른 질문에 답한다.} RU는 ``이 직업 종사자 1인(1시간)당 AI를 전국 평균의 몇 배 쓰는가'', 점유율은 ``AI 사용 총량이 어느 직업에 있는가''다. 노출도처럼 1인당 강도 지표와 비교하거나 직업별 생산성 효과를 볼 때는 RU가, 일자리 영향 총량을 직업에 나눌 때는 점유율이 맞는 지표다.
  \item \textbf{RU의 직업 간 차이는 상당 부분 분모의 차이다.} 이 자료에서 대화 점유율은 취업자 수와 무관하다(규모 탄력성 $-0.01$). 그래서 RU1은 사실상 ``A $\div$ 취업자 비중''처럼 움직이고, 취업자가 적은 직업에서 커진다. 이는 연계표의 균등 배분(paper 21)과 AEI의 과업 기준 직업 정의에서 나오는 성질이다. 취업자 가중 배분으로 바꾸면 탄력성이 0보다 커지고 RU와 점유율의 격차도 줄어들 것이다. 그 민감도 분석은 아직 하지 않았다.
  \item \textbf{점유율도 이 성질을 물려받는다.} 점유율 A가 규모와 무관하다는 것은 ``큰 직업에 AI 사용이 더 많다''는 증거가 없다는 뜻이다. 동시에 점유율 지표가 직업 규모를 제대로 반영하지 못할 수 있다는 뜻이기도 하다. 521 매장 판매 같은 큰 직업의 A가 취업자 비중(129개 안 6.9\%)보다 훨씬 작은 0.83\%인 것은 실제 사용이 적어서일 수도 있고, 배분이 규모를 따르지 않아서일 수도 있다.
  \item \textbf{표본이 작은 직업.} 표본 30 미만 소분류 6개는 분모가 불안정하다. RU 상위 목록을 쓸 때는 $\dagger$ 직업을 따로 표시하거나 빼고 보는 것이 안전하다. 점유율 지표는 분모가 없어 이 문제가 없다.
  \item \textbf{상대값.} RU는 전국 평균 대비 배수라 수준(몇 \% 절감)을 말하지 않는다. 예전 활용도 2는 100\%를 넘는 소분류가 있어 순위 비교용이었는데(paper 25), RU로 바꾸면서 이 점이 정의에 드러났다. 이 노트도 순위만 비교했다.
\end{enumerate}

\section{산출 파일}

\begin{itemize}[leftmargin=*]
  \item 코드: \texttt{code/04\_15\_analysis\_kor\_use\_vs\_share.do}(Stata 17). 입력은 \texttt{kor\_ai\_usage\_ksco3}(\texttt{03\_05\_merge\_kor\_ksco3\_ai\_usage.do}). 로그 항등식, 분산 분해와 회귀 기울기의 관계를 assert로 확인한다.
  \item 표: \texttt{results/table/32\_kor\_use\_vs\_share.xlsx}. 시트 \texttt{ksco3}(129개 소분류의 RU1$\cdot$RU2(열 이름 \texttt{use1}, \texttt{use2}), A, B, 취업자, 총 근로시간, 순위, $G_1$, $G_2$), \texttt{stats}, \texttt{top}($G_1$, $G_2$ 상위$\cdot$하위 15), \texttt{topuse}(RU1$\cdot$RU2 상위 20), \texttt{ksco1}.
  \item 이 노트의 표: \texttt{results/table/32\_summary.tex}, \texttt{32\_g1\_pos.tex}, \texttt{32\_g1\_neg.tex}, \texttt{32\_g2\_pos.tex}, \texttt{32\_g2\_neg.tex}, \texttt{32\_top\_use1.tex}, \texttt{32\_top\_use2.tex}, \texttt{32\_ksco1.tex}.
\end{itemize}

\end{document}
